% !TeX root = main.tex
\section{引言}\label{sec:intro}

能源系统运行、生产调度和交通管理等任务需要同时考虑效率、成本与能耗等目标。节能驾驶中速度与能耗的权衡、基础设施维护中成本与失效风险的权衡，都是这类问题的具体实例\cite{he2023,remmerden2025}。这些目标之间存在冲突，一个目标的改善往往以其他目标变差为代价。这类任务还是序贯的：当前动作除了产生即时结果，还会改变系统状态，从而影响后续可采取的动作及其后果。例如，在具有储能或热惯性的能源系统中，降低当前运行功率可以减少即时能耗，却可能改变系统后续状态，增加之后恢复或补偿运行的成本。策略的优劣因此要依据一系列动作产生的长期结果来判断。多目标序贯决策要回答的，就是在系统持续演化的过程中如何选择策略，从而在多个相互冲突的目标之间取得权衡。

对于一组互不支配的候选策略，Pareto 最优性只说明它们之间不存在支配关系，并不能确定选哪一个；实际决策还需要决策者对不同目标的取舍要求，即偏好。权重、效用函数、参考点、满意水平和目标优先级等形式，可以从不同角度表达这种取舍\cite{wang2017}。本文研究偏好已经给定的情形，即决策者的取舍依据是决策问题的已知输入，讨论如何把它纳入多目标序贯决策的建模、求解与评价。

偏好给定并不意味着问题已经化为单目标问题。首先要明确的是：给定偏好定义了怎样的长期决策准则，这一准则又具有怎样的序贯结构。固定线性权重下，可以在一定条件下把多维回报标量化，再用单目标序贯决策方法求解；对于更一般的偏好，长期准则未必具有阶段可加性，当前环境状态也可能不足以支持后续决策，需要引入累计回报、扩展状态、回报分布等专门的表示与求解机制。具体求解研究已经说明了这一点：累计回报上的非线性效用推动了保留历史结果的搜索和回报分布学习\cite{hayes2023,reymond2023}，严格优先关系要求策略更新始终维持目标之间的比较次序\cite{skalse2022}。信息条件又进一步区分了求解路线：模型已知时可以直接规划，允许交互时可以补充样本，只有固定历史数据时还要处理数据对不同权衡的支持程度，从多奖励批量价值学习\cite{lizotte2012}到偏好与回报条件化的离线策略\cite{zhu2023}，数据覆盖始终是这类方法的适用边界。

已有综述从不同角度整理了这一领域。Roijers 等从三种无法、难以或不宜把多目标问题转化为单目标问题的场景出发，按场景、标量化函数的性质和所考虑的策略类型建立分类，据此确定最优解应是单个策略、凸包还是 Pareto 前沿，并按这一分类综述规划与学习方法\cite{roijers2013}。Hayes 等面向希望采用多目标视角的研究者和实践者，提出以用户效用为中心的设计流程：先收集效用信息、确定允许的策略类型，再推出解的概念并选择算法；该文讨论了 SER 与 ESR 的选择依据和策略类型的影响，并对算法和评价指标作了简要综述\cite{hayes2022}。Wang 等的近期综述从理论基础、算法框架、评价基准和应用几个方面总结多目标强化学习研究，按模型可用性、策略范围和偏好表示对算法分类，涵盖基于模型的规划、无模型学习和交互式偏好获取\cite{wang2026}。

现有综述主要围绕"需要求出什么样的解"和"有哪些算法"来组织文献。对于偏好已经给定的问题，仍有三个方面值得专门梳理。其一，给定偏好形成的长期准则不只是 SER 与 ESR 的选择，还包括评价对象、时间聚合方式、优先关系和可接受性要求，它们共同决定了求解时必须保留哪些信息。其二，模型可用、允许交互和只有固定历史数据三类信息条件下，同一准则的求解方法和理论保证差别很大；上述综述的算法分类没有把固定历史数据单列为一类条件。其三，策略库和条件策略等多策略方法在给定偏好问题中的作用是复用计算和数据，对它们的评价最终要落实到每个给定偏好下实际选出的策略。

围绕这三个方面，本文做了以下工作：
\begin{enumerate}
\item 提出“长期准则---信息表示---策略类别---模型与数据条件”四维分析框架，把偏好表达、评价对象、时间聚合和随机性处理统一为长期决策准则，并据此分析各类准则的可分解性、状态充分性、递推结构和所需策略类别。
\item 按模型可用、允许交互和固定历史数据三类信息条件组织求解方法，每类条件下再按准则类型比较方法的机制与保证范围，并将跨偏好方法单独作为复用机制讨论。
\item 从原准则出发讨论单策略、策略集合与条件策略的评价，以及离线评价、实验规范和应用证据的层次，并据此提出四个研究方向。
\end{enumerate}

全文共八章。第~\ref{sec:framework}~章先界定研究范围并给出分析框架；第~\ref{sec:criteria}~章讨论给定偏好如何形成长期决策准则，第~\ref{sec:structure}~章在此基础上分析准则的序贯结构与求解条件。随后，第~\ref{sec:solving}~章按三类信息条件比较求解路线，第~\ref{sec:evaluation}~章回到准则本身，讨论策略评价与应用证据。第~\ref{sec:future}~章提出未来研究方向，最后一章总结全文。
